% ======================================================================
% Methods: defect statistics, charge profiles and edge/bulk criteria
% Generated from the analysis scripts in runs/cavity and runs/cassini.
% All facts checked against the code (2026-10-05). Where code and
% RESEARCH.md disagree, the code is followed and a NOTE is left.
% ======================================================================

\subsection{Defect data and charge audit}
\label{sec:methods-audit}

\paragraph{Raw data.}
All defect analyses start from the solver's defect file \texttt{defects.dat}
(read transparently from \texttt{defects.dat.gz} when the run was compressed).
Each frame consists of a header line with two numbers, the time $t$ and the number
$n$ of defects, followed by $n$ lines with four columns: the position $(x,y)$ in
cell units, the topological charge $q$ and the orientation angle $\psi$ of the defect.
The tracker only reports clusters with $|q|\simeq 1/2$; the charge column holds
values within $10^{-2}$ of $\pm 1/2$ (in a typical run $\approx 0.1\%$ of entries
differ from $\pm0.5$ by $10^{-3}$--$5\times10^{-3}$). Defect frames are written every
\texttt{defectsOut} $=20$ MPCD steps, i.e.\ every $\Delta t_{\rm f}=2$ time units
($\delta t = 0.1$); the main circle sweep records $10^4$ time units of production
($5001$ frames) after a $10^4$-step warm-up that is not written.
Non-numeric lines (the banner and column headers) are skipped; a truncated final frame
(fewer than $n$ defect lines) is discarded (\texttt{defect\_audit.read\_frames}).
Positions are shifted to the cavity centre $\mathbf{Q}$ given in \texttt{input.json}.

\paragraph{Transient.}
Unless stated otherwise every script discards the first $20\%$ of the recorded frames
(\texttt{--tfrac 0.2}), i.e.\ $2000$ time units in the main sweep, and uses the remaining
$\approx 4000$ frames. \texttt{defect\_spatial.load} and all scripts built on it drop
the first $\lfloor 0.2\,N_{\rm f}\rfloor$ frames by count; \texttt{layer\_analysis.load}
drops frames with $t < t_{\min}+0.2\,(t_{\max}-t_{\min})$, which is equivalent for
uniform sampling. \texttt{circulation.py} uses all production frames (\texttt{--tmin 0};
the warm-up is not in the files).

\paragraph{Charge audit (\texttt{defect\_audit.py}).}
For imposed winding $k$ (sign $\sigma_k=\mathrm{sgn}\,k$) every frame is checked for
topological consistency using only the signs of the tracked charges. With
$N_{\rm same}$ the number of defects with the sign of $k$ and $N_{\rm opp}$ the number of the
opposite sign,
\begin{equation}
  q_{\rm tot} = \tfrac12\,\sigma_k\,(N_{\rm same}-N_{\rm opp}),
\end{equation}
and the frame is classified by three flags:
\begin{itemize}
  \item \emph{net charge}: $|q_{\rm tot}-k|<1/4$ (charge conservation, total charge $=k$);
  \item \emph{imposed}: $N_{\rm same}\ge 2|k|$ (all boundary-imposed $\pm1/2$ present);
  \item \emph{pairs}: $N_{\rm same}-2|k| = N_{\rm opp}$ (the dynamical defects,
        $N_{\rm dyn}=N_{\rm same}+N_{\rm opp}-2|k|$, come in $\pm1/2$ pairs).
\end{itemize}
For $\pm1/2$ charges the last flag is equivalent to the first; they are reported
separately because the failure mode is diagnostic: a frame with $N_{\rm same}<2|k|$
has lost wall defects (``missing imposed''), a frame that passes ``imposed'' but not
``pairs'' has a missed, spurious or split pair member (``unpaired''). Per run we report
the fraction of frames with $q_{\rm tot}=k$, the fractions ``missing imposed'' and
``unpaired'', the fraction with odd $N_{\rm dyn}$, $\langle q_{\rm tot}\rangle$ and the mean
number of dynamically created pairs $\langle N_{\rm pairs}\rangle=\langle N_{\rm opp}\rangle$.
The audit needs no geometry and therefore applies unchanged to ovals and ellipses.
In passive validation runs it is satisfied in $100\%$ of frames; in the active fixed-$R$ scans
at $\zeta=0.07$ it holds in $\ge98\%$ of frames. The anchoring itself fails (audit
$\le 30\%$) for imposed linear densities $\rho\gtrsim0.25$, which bounds the parameter
space analysed ($\rho\le 0.15$).
% NOTE: the audit uses sign(q) only; quantities built on the raw charge column (Q(c), site-rule q_b)
% carry the small (~1e-3) deviations of individual charges from +-1/2 (see the P1 note below).

\subsection{Spatial distributions}
\label{sec:methods-spatial}

\paragraph{Time-averaged maps (\texttt{defect\_spatial.py}).}
Positions of $+1/2$ and $-1/2$ defects from all retained frames are binned on a square
grid of $0.5\times0.5$ cells covering $[-R-0.5,R+0.5]^2$; counts are divided by the number of
frames and the bin area, giving the mean number density $n_\pm(\mathbf{x})$ (defects per
cell$^2$ per frame). The maps are shown on a logarithmic scale spanning three decades.

\paragraph{Radial densities in circles.}
With $s=R-r$ the distance from the wall, defects are binned in $s$ with width
$\Delta s=0.25$ (\texttt{--sbin}); the counts are divided by the exact annulus area
$\pi(r_{\rm out}^2-r_{\rm in}^2)$ and by the number of frames, giving $n_\pm(s)$.

\paragraph{Defect speeds.}
For each pair of consecutive frames and each sign, every defect is matched to the nearest
same-sign defect of the next frame; the match is accepted if the distance is below
$r_{\rm link}=3$ cells, and the speed is the displacement divided by $\Delta t_{\rm f}$. Speeds are
binned in $s$ (same bins as $n_\pm$) and averaged over bins with at least $20$ samples.
This matching is nearest-neighbour (not one-to-one).
% NOTE: with Delta t_f = 2 and bulk +1/2 speeds ~0.8-1 cells/time, r_link = 3 rejects the fastest
% steps (RESEARCH 2026-10-01 18:20: ~90% of track breaks at r_link = 3 were genuine 3-6 cell jumps), so
% the bulk speed profile is biased low; the tracking used for the mobility criterion uses r_link = 8.
Layer $+1/2$ defects move at $0.2$--$0.35$ cells per time unit, indistinguishable from passive
pinned defects ($\approx0.3$, the tracker's positional jitter), whereas bulk $+1/2$ move at
$0.8$--$1.0$.

\paragraph{Angular sites in circles.}
As a first count of pinned wall defects, the polar angles of all $+1/2$ with $s<4.2$ are
histogrammed in $360$ bins, smoothed periodically with a Gaussian of width $2$ bins
(13-tap kernel) and the number of runs (contiguous intervals) of bins above $15\%$ of the
maximum is counted (\texttt{defect\_spatial.layer\_sites}). This circle-specific count is
superseded by the shape-independent site rule of Sec.~\ref{sec:methods-sites}.

\paragraph{Profiles along Cassini and ellipse levels (\texttt{lobe\_profiles.py}).}
For Cassini ovals with foci at $(\pm a,0)$ and wall $d_1d_2=b^2$ we use the level coordinate
\begin{equation}
  c(\mathbf{x}) = \sqrt{d_1 d_2},\qquad d_{1,2}=|\mathbf{x}\mp a\hat{\mathbf{e}}_x| ,
\end{equation}
so that $c=0$ at the foci, $c=a$ on the lemniscate (for $c<a$ the level consists of two
separate lobe ovals) and $c=b$ on the wall; for $a=0$ (circles) $c=r$. For ellipses with
semi-axes $\varepsilon B$ (along $x$) and $B$ (along $y$) the levels are scaled copies of the
wall,
\begin{equation}
  c(\mathbf{x}) = \sqrt{(x/\varepsilon)^2+y^2},\qquad c=b\equiv B \text{ on the wall},
\end{equation}
with $\varepsilon$ and $B$ read from the wall parameters ($\varepsilon = a^{-1}_x/a^{-1}_y$,
$B=R\,a^{-1}_y$ for \texttt{aInv} $=(a^{-1}_x,a^{-1}_y,\cdot)$).
The interval $c\in[0,b]$ is divided into $N_{\rm sh}$ shells uniform in $c/b$
($N_{\rm sh}=200$ for all quantitative results, i.e.\ $\Delta(c/b)=0.005$; the plotting default is
60). Shell areas are exact for ellipses, $A_j=\pi\varepsilon(c_{j+1}^2-c_j^2)$, and for
Cassini ovals are obtained by counting the points of a $0.1$-cell grid whose $c$ falls in the
shell. Densities are $n_\pm(c)=$ (defects in the shell per frame)$/A_j$. Profiles per lobe
are obtained by restricting to $x<0$ or $x>0$ (meaningful for $c<a$ only). For Cassini ovals
the physical width of a shell varies along the level ($\propto 1/|\nabla c|$, widest at the
neck), which smears a thin wall layer; conclusions were therefore cross-checked against
profiles in the true wall distance (Sec.~\ref{sec:methods-walldist}).

\subsection{Accumulated charge and level coordinates}
\label{sec:methods-Q}

\paragraph{Accumulated charge $Q(c)$.}
The central observable is the mean net charge enclosed by a level,
\begin{equation}
  Q(c_j) = \Big\langle \sum_{i:\,c(\mathbf{x}_i)<c_j} q_i \Big\rangle_t ,
\end{equation}
computed as the cumulative sum over shells of the charge histogram (weights $q_i$, the raw
charge column) divided by the number of frames, and evaluated at the \emph{outer} shell
edges $c_j/b = j/N_{\rm sh}$ (\texttt{lobe\_profiles.profile}). In circles $Q(r)$ is the
charge inside radius $r$. By construction $Q(b)=k$ (up to defects reported outside the
wall, which are not counted). Per-lobe charges $Q_{L,R}(c)$ use $x<0$ or $x>0$.
For circles, $Q(r)$ shows a negative dip towards the centre ($-1/2$ excess), a plateau
at $+1$ just inside the wall, and a jump to $k$ in the last few percent of $R$ (the layer).

\paragraph{Fixed-distance enclosed charge $Q(s)$.}
As a cut-free alternative in general shapes, \texttt{cassini\_charge.py} reports the mean
net charge of defects farther than $s$ from the wall, for $s\in\{1,2,3,4,5,6,8,10\}$
cells, with $s$ the first-order wall distance of Sec.~\ref{sec:methods-walldist}.
In circles at $\zeta=0.07$ this is flat ($Q\approx1.0$) for $s=3$--$5$; in elongated ovals it
falls steadily without a plateau, so a single ``bulk'' value at fixed $s$ is cut-dependent there.

\paragraph{Wall distance in general shapes.}
\label{sec:methods-walldist}
For the Cassini family ($a=0$: circle) the wall function is $W=d_1d_2-b^2$, and the distance
from the wall and the outward normal angle are taken to first order,
\begin{equation}
  s(\mathbf{x}) = -\frac{W}{|\nabla W|},\qquad
  \nabla W = \Big(\tfrac{d_2}{d_1}(x-a)+\tfrac{d_1}{d_2}(x+a),\;
  \big(\tfrac{d_2}{d_1}+\tfrac{d_1}{d_2}\big)y\Big),\qquad
  \phi_n = \operatorname{atan2}(\partial_yW,\partial_xW)
\end{equation}
(\texttt{edge\_criteria.wall}, \texttt{cassini\_charge.wall\_distance}).
% NOTE: s = -W/|grad W| is exact only at the wall. For a circle it gives s = (R^2 - r^2)/(2r)
% = (R - r)(R + r)/(2r), not R - r: at R = 31.8 the cut s < 6 corresponds to a true distance of
% ~5.4 cells, and s diverges at the centre. Ellipse runs (shape 0, aInv_x != aInv_y) are treated as
% circles of radius R by these two functions (only cassiniA / shape 1 is read).
For the mean wall distance along a level curve (Sec.~\ref{sec:methods-dstar}) the exact
Euclidean distance to a polygon of $4000$ wall points is used, the wall being
$r^2(\theta)=a^2\cos2\theta+\sqrt{b^4-a^4\sin^2 2\theta}$ (single oval, $b>a$).

\subsection{Bulk charge from the shape of $Q(c)$: knee, $d^*$ and $f^*$}
\label{sec:methods-knee}

\paragraph{Knee of $Q(c)$ (\texttt{charge\_knee.py}).}
The bulk charge is read from the plateau of $Q$ just inside the wall, located as the
``knee'', i.e.\ the outermost point where $Q$ is clearly flatter than on both sides.
With $x_j=c_j/b$ the outer shell edges ($N_{\rm sh}=200$):
\begin{enumerate}
  \item \emph{Smoothing.} $Q$ is padded with $9$ copies of its first and last values and
        convolved with a $9$-bin box kernel, $\tilde Q = Q * \mathbb{1}_9/9$, i.e.\ smoothed over
        $0.045$ in $c/b$. The slope is $S=\mathrm{d}\tilde Q/\mathrm{d}x$ (centred finite differences,
        \texttt{numpy.gradient}).
  \item \emph{Search window.} From $i_0+2$ (with $i_0=\arg\min\tilde Q$, the $-1/2$ core; at least
        index 3) up to the index of $x=0.985$ minus 3.
  \item \emph{Candidates.} Local minima of $S$ ($S_i\le S_{i\pm1}$) that satisfy all of
        \begin{align}
          S_i &\le 0.7\, \max_{x_i-0.12\le x_j< x_i} S_j &&\text{(flatter than the inner side)},\\
          S_i &\le 0.2\, \max_{x_i< x_j\le 0.985} S_j &&\text{(flatter than the jump to the wall)},\\
          \tilde Q_i - \tilde Q(x_i-0.2) &\ge 0.5 &&\text{(a plateau is reached by a rise)} .
        \end{align}
  \item \emph{Selection.} Among the candidates the one closest to the wall (largest $x_i$) is the knee.
\end{enumerate}
The bulk charge is $Q_{\rm knee}=\tilde Q(x_{\rm knee})$, the layer charge $k-Q_{\rm knee}$; if no
candidate exists the run has \emph{no knee} (no layer/bulk split). The plateau contrast
$\max_{\rm wall\,side}S/S_{\rm knee}$ is also recorded. The rise condition excludes noise minima in
a flat interior; the $0.985$ limit excludes the final jump onto the wall.
The smoothing width was changed from $5$ to $9$ bins after the main-sweep analysis: with $5$ bins
the finder picked up the flat shoulders of the cell-quantised layer steps (width $\approx0.02$ in
$c/b$), e.g.\ a spurious knee at $Q\approx1.9$ for $k=5$, $R\approx0.8$--$0.9\,\la$ whose true plateau is
at $+1.0$; $9$ bins ($0.045$) removes these steps but preserves a genuine plateau
($\approx0.1$ wide). All runs were re-analysed with $9$ bins; the established cases were unchanged.
The same function is applied to ovals and ellipses in their own level coordinate.

\paragraph{Mirror runs ($k<0$): peak of $Q$.}
For negative winding $Q$ rises to a positive maximum and then plunges to $-|k|$ at the wall, so the
bulk charge is taken as the maximum of the smoothed $\tilde Q$ (same $9$-bin smoothing) over
$x<0.97$ (\texttt{criteria\_all.py}). It is accepted as a plateau only if (i) the band of $x<0.97$
where $\tilde Q>\tilde Q_{\max}-0.1$ is at least $0.04$ wide in $r/R$, (ii) the steepest
$|\mathrm{d}\tilde Q/\mathrm{d}x|$ between the peak and the wall is at least $5\times$ the mean
$|\mathrm{d}\tilde Q/\mathrm{d}x|$ over the band, and (iii) $\tilde Q_{\max}>0.05$. The earlier survey
script \texttt{mirror\_analysis.py} reports the unsmoothed peak of $Q$ for $x<0.97$, its position and
the width of the band within $0.1$ of it, together with the mean numbers of $\pm1/2$ within $s<3$
(wall) and $s\ge3$ (bulk) per frame.

\paragraph{Layer depth $d^*$ (\texttt{layer\_depth.py}).}
\label{sec:methods-dstar}
The level $c^*$ where $Q$ crosses $+1$ on its way to $k$ is located by scanning inward from the
wall: with $i$ the outermost shell edge where $Q<1$,
$c^*/b = x_i + (1-Q_i)(x_{i+1}-x_i)/(Q_{i+1}-Q_i)$ (linear interpolation; unsmoothed $Q$).
The layer depth is the mean wall distance along the level $c=c^*$. In circles $d^*=R-c^*$
(as used in \texttt{main\_sweep\_analysis.py}). For ovals the level is sampled on a square grid of
spacing $h=0.15$: all grid points with $|c(\mathbf{x})-c^*|<\tfrac12 h\max(1,\overline{|\partial_x c|})$
(median over nearby points) are retained and their exact distances to the wall polygon are averaged.
% NOTE: the docstring/RESEARCH call this an "arclength-weighted" mean. Grid points in a band of fixed
% width in c are weighted by 1/|grad c| per unit arclength (denser where levels are farther apart,
% i.e. at the neck), so the mean is arclength-weighted only where |grad c| is constant along the level
% (exactly so for circles).
$d^*$ is reported in cells, in units of $\la$ ($\la=22.21$) and relative to the cavity size.
A plateau width $w$ (the wall-distance difference between the inner and outer ends of the contiguous
band around $c^*$ with $|Q-1|<0.15$) is also computed, but is unreliable for ovals (it can become
negative when $Q$ merely passes through $1$) and is not used.

\paragraph{Layer area fraction $f^*$ and effective width (\texttt{layer\_area.py}).}
With $c^*$ as above,
\begin{equation}
  f^* = \frac{A(c^*<c<b)}{A(c<b)},\qquad w_{\rm eff}=\frac{A(c^*<c<b)}{P},
\end{equation}
with areas counted on a $0.1$-cell grid and $P$ the perimeter of the $4000$-point wall polygon.
In circles $f^*=1-(c^*/R)^2$ (\texttt{main\_sweep\_analysis.py}, \texttt{criteria\_all.py}).
$f^*=1/2$ (the ``layer'' as large as the region it encloses) is used as the bound beyond which the
outer region is not a layer: clean circles have $f^*\le0.32$. For a true layer $f^*\propto$
width$/$size decreases with cavity size, whereas a charge distribution spread through the cavity
gives a size-independent $f^*$; the size scaling of $f^*$ thus distinguishes the two.
% NOTE: layer_area.py / layer_depth.py read the geometry with lobe_profiles.geometry, which ignores
% the ellipse parameters (a = 0, b = R). Because ellipse levels are self-similar, f* = 1 - (c*/b)^2 is
% still correct for ellipses, but w_eff and d* would be computed for a circle of radius R; only f* is
% quoted for ellipses. layer_area.py has no guard for runs where Q never drops below 1.
In mirror runs the $+1$ crossing is taken from the centre outward (first $x<0.97$ with $Q\ge1$).

\subsection{Edge/bulk classification of individual defects}
\label{sec:methods-criteria}

All criteria below assign each defect of a frame to the wall layer (``edge'') or to the bulk; the bulk
charge is the mean net charge of bulk defects, $q_b=\langle\sum_{i\in\rm bulk}q_i\rangle_t$, and
$P_1$ is the fraction of frames with $q_b=+1$ (tested as $|q_b-1|<1/4$ unless noted).
In the per-defect rules (site, density, mobility, tilt) only defects of the sign of $k$ (the wall-imposed sign) are
ever assigned to the edge; the charge-profile criteria (knee, fixed cut, $q_{\rm dens}$) assign everything outside a level to the edge.

\paragraph{Fixed geometric cuts.}
(i) $q_{\rm fixed}=Q(r=\tfrac{11}{12}R)$, linearly interpolated in the $200$-shell $Q(r)$
(\texttt{criteria\_all.py}). (ii) In \texttt{circulation.py} bulk defects are those with
$r<R-2.5$; $q_b(t)=\tfrac12\sum_{\rm bulk}\mathrm{sgn}\,q_i$. Fixed cuts fail when the layer thickness
does not scale with $R$: below $\sim1\,\la$ the layer is thicker than $R/12$ (then $q_{\rm fixed}\gg1$),
and for $k<0$ at $R\lesssim1.5\,\la$ it is strongly negative; they are reported only for comparison.
% NOTE: RESEARCH (2026-09-30) notes that the fixed --wall 2.5 in circulation.py underestimates the layer.

\paragraph{Passive edge $d_p$ (\texttt{layer\_analysis.py}).}
The layer is taken from the passive run with the same $(R,k)$ (located by the same \texttt{rho}/\texttt{k}
path, seed \texttt{s01}): $d_p$ is the $95$th percentile of $s=R-r$ over all same-sign defects of the passive
run (after the 20\% transient), and a layer is said to exist if at least half of them lie within $5$ cells
(\texttt{--layer-max}) of the wall. Bulk: $s>d_p$. This cut was abandoned: the passive layer shrinks
from $8.3$ to $2.3$ cells with $k$ and is thinner than the active one at $k\gtrsim6$ (and too wide at $k=2$),
giving $q_b$ from $0.44$ to $1.20$ in a set where the density cut gives $1.00\pm0.02$ for every $k$.

\paragraph{Density-minimum edge $d_\rho$.}
Two implementations are used.
\begin{itemize}
  \item \texttt{layer\_analysis.density\_edge}: the radial density of same-sign defects is computed in
        bins of $0.5$ cells in $s\in[0,\min(R,12)]$ (per annulus area and frame). The wall peak is the
        maximum among bins with centre $s<3$; $d_\rho$ is the bin centre of the minimum of the $3$-bin box-smoothed
        density over bins with $s_{\rm peak}<s\le\min(8,R/2)$. Bulk: $s>d_\rho$; layer: same-sign defects with
        $s\le d_\rho$. From the layer we also record the occupancy $N_L(t)$, its relative fluctuation
        $\mathrm{std}/\mathrm{mean}$ and the radial width (standard deviation of $s$ in the layer).
  \item \texttt{criteria\_all.py} ($q_{\rm dens}$), on the $200$-shell profile in $x=r/R$: the wall-sign density
        $n_{\sigma_k}$ is smoothed with a $5$-bin box ($\Delta x=0.025$); the wall peak is its maximum for
        $x>1-6/R$ (within $6$ cells of the wall); the gap is the minimum of the smoothed density between
        $x_{\rm lo}=\max(1-10/R,\,0.3)$ (within $10$ cells of the wall) and the peak, required to lie at
        least $2$ bins inside it. Then $q_{\rm dens}=Q(x_{\min})$ (charge inside the minimum, original signs).
        Undefined if the peak is not more than $2$ bins outside $x_{\rm lo}$.
\end{itemize}
In circles at $\zeta=0.07$ the $+1/2$ density has a sharp wall peak at $s\approx1$--$2$, a gap at
$s\approx3$--$5$ about $5\times$ below the bulk plateau, and the speed step sits at the same $s$, so
$d_\rho\approx4$ cells. The gap does not scale with $\la$ (layer width is set in cells, bulk density
$\propto\la^{-2}$): it is shallow at $\zeta=0.06$ and almost absent at $0.11$, where the split becomes
ill-defined. The density criterion gives $q_b\to0$ in very small cavities ($R\lesssim0.6$--$0.8\,\la$,
no bulk defects).

\paragraph{Pinned-site rule (\texttt{edge\_sites.classify}).}
\label{sec:methods-sites}
Wall defects are stationary, so they leave sharp spots in the time-averaged density while mobile bulk
defects average out.
\begin{enumerate}
  \item \emph{Sites.} The $+1/2$ positions of all retained frames are binned on a $0.5$-cell grid,
        normalised to defects per cell$^2$ per frame and smoothed with a Gaussian of
        $\sigma=1$ cell. Sites are local maxima (equal to the maximum in a
        $9\times9$-bin, i.e.\ $4.5\times4.5$-cell, window; window size $2\,\mathrm{round}(2\sigma/0.5)+1$)
        that exceed $0.02$ defects per cell$^2$ per frame. No radial cut or circle geometry is used.
  \item \emph{Capture radius.} For every $+1/2$ occurrence the distance to the nearest site is computed;
        its histogram (bins of $0.25$ cells on $[0,12]$) is smoothed with a $5$-bin box and
        $d_{\rm site}$ is the position of its minimum between $1.5$ and $8$ cells (fallback $3$ cells).
        A fixed $d_{\rm site}=2$--$3$ is too tight (in a passive $k=5$ run it captures only $7.8$--$9.6$
        of $10$ pinned defects).
  \item \emph{Cap.} $d_{\rm site}$ is capped at $0.4\times$ the median nearest-neighbour distance
        between sites (\texttt{--dfrac 0.4}). Without the cap, when the bulk $+1/2$ ring presses against
        the layer (no dip in the histogram, e.g.\ $\zeta=0.11$, $k\ge7$) the minimum runs to the $8$-cell
        window edge and the rule captures bulk defects ($q_b\approx-0.5$).
        % NOTE: RESEARCH says the cap "never binds" at zeta = 0.07 (max |Delta q_b| = 0), but also
        % that the uncapped d_site ~ half the site spacing in the R-scan; 0.4 x spacing < 0.5 x spacing, so
        % the cap may bind there without changing q_b. The code records a 'capped' flag per run.
  \item \emph{Classification.} A $+1/2$ within $d_{\rm site}$ of a site is edge; all other defects
        (and every $-1/2$) are bulk. Per frame $q_b(t)=\sum_{\rm bulk}q_i$ and the edge occupancy $N_E(t)$.
\end{enumerate}
Reported per run: number of sites (compare with $2k-2$), mean wall distance of the sites, $d_{\rm site}$,
$\langle N_E\rangle$, $\langle q_b\rangle$ with a block error ($10$ consecutive blocks, standard error of
the block means), $P_1$, and the fraction of $+1/2$ occurrences classified identically by the site rule
and by a density cut at $s=4.2$ (circles: $96$--$99.7\%$ at $\zeta=0.07$).
% NOTE: in edge_sites.classify P1 is computed as np.isclose(q_b, 1) on the raw charge column, whose
% entries differ from 0.5 by up to 5e-3 in ~0.1% of defects; frames containing such a defect fail the
% test, so P1(site) is biased low by roughly the fraction of such frames (percent level). Other P1's
% use sign(q) and |q_b - 1| < 1/4.
For mirror runs ($k<0$) the frames are sign-flipped before classification, so the sites are those of the
wall $-1/2$, and $q_{\rm site}=\sigma_k\,q_b$ is returned in the original charges (\texttt{criteria\_all.py}).
For Cassini ovals the bulk charge is also split into left and right halves ($x<0$, $x>0$)
(\texttt{cassini\_charge.py}). Limitations: in passive runs at $k=2,3$ the pinned defects sit away from the
wall (the rule finds pinned, not wall, defects; $q_b=0$ is still correct), and at $\zeta=0.05$, $k=9,10$
(flow reversals) spurious interior maxima give $16$--$27$ sites; such runs are flagged.

\paragraph{Linking defects into tracks.}
Three linkers are used.
\begin{itemize}
  \item \texttt{edge\_criteria.link} ($+1/2$ only): consecutive frames are matched by an optimal
        one-to-one assignment (Hungarian algorithm, cost = distance, pairs beyond $r_{\rm link}$ forbidden);
        unmatched defects start new tracks; no gap closing.
        $r_{\rm link}=8$ cells is used for all reported criteria (the command-line default is $3$).
        With $\Delta t_{\rm f}=2$ and cell-quantised positions, $r_{\rm link}=3$ breaks most bulk tracks
        (median track length one frame); $r_{\rm link}=8$ reduces the number of tracks $\sim9\times$ and
        $>90\%$ of track starts/ends are then paired with an opposite-sign start/end.
  \item \texttt{layer\_analysis.link} (both signs): optimal assignment per frame with gap closing:
        a track may miss up to $3$ frames, the admissible jump grows as $r_{\rm link}\sqrt{g}$ for a gap of $g$
        frames, opposite signs are never linked, and a penalty $0.5(g-1)$ is added to the cost to prefer
        continuing tracks seen in the last frame ($r_{\rm link}=3$; runs with $\Delta t_{\rm f}>2.5$ are flagged as
        unreliable for linking).
  \item \texttt{defect\_spatial} speeds: nearest neighbour, see Sec.~\ref{sec:methods-spatial}.
\end{itemize}

\paragraph{Mobility criterion (\texttt{edge\_criteria.py}).}
For every $+1/2$ occurrence at frame $n$ on track $j$, the net displacement over a window $\tau$,
\begin{equation}
  D = |\mathbf{r}_j(n+h)-\mathbf{r}_j(n-h)|,\qquad h=\max\!\big(1,\mathrm{round}(\tau/2\Delta t_{\rm f})\big),
\end{equation}
with $\tau=20$ time units ($h=5$ frames). Occurrences whose track does not exist at both $n-h$ and $n+h$,
and frames within $h$ of either end of the series, are skipped. Wall defects in circles have a median
$D$ of $1.0$--$1.4$ cells against $10$--$11$ for bulk $+1/2$.
In \texttt{criteria\_all.py} a defect is ``slow'' (edge) if $D<5$ cells \emph{and} $s<6$ (first-order
wall distance). The bulk charge is then obtained from the total charge,
\begin{equation}
  q_{\rm slow} = k - \tfrac12\,\frac{N_{\rm slow}}{N_{\rm f}'},
\end{equation}
with $N_{\rm slow}$ the number of slow occurrences and $N_{\rm f}'=N_{\rm f}-2h$ the number of frames
used. Normalising by frames (not by the number of tracked $+1/2$) avoids a bias towards the edge,
because bulk tracks break more often than wall tracks.
% NOTE: occurrences skipped because their track breaks inside the window are not counted as slow,
% so N_slow/N_f' is a lower bound on the edge count (q_slow biased high) when wall tracks break;
% this assumes total charge k (audit) rather than counting bulk defects directly.
For mirror runs the sign of all charges is flipped before linking so that the wall $-1/2$ are followed,
and $q_{\rm slow}=k+\tfrac12 N_{\rm slow}/N_{\rm f}'$.

\paragraph{Tilt criterion.}
The orientation of a wall $+1/2$ relative to the local outward wall normal is
$\Delta=\mathrm{wrap}_{(-\pi,\pi]}(\psi-\phi_n)$, with $\psi$ the angle column of \texttt{defects.dat} and
$\phi_n$ the normal angle of Sec.~\ref{sec:methods-walldist}. In circles wall $+1/2$ are tilted by
$|\Delta|\approx1.1$~rad ($\approx63^\circ$) with the sign of the vortex sense, nearly independent of $k$,
while bulk $+1/2$ cluster near $|\Delta|\approx\pi-0.5\approx2.4$~rad. The combined criterion
(``slow and tilted'', $+k$ only) counts an occurrence as edge if $D<5$, $s<6$ and $|\Delta|<1.8$~rad,
$q_{\rm tilt}=k-\tfrac12 N_{\rm tilt}/N_{\rm f}'$. Only an upper cut is used because the wall tilt
decreases from $\approx1.1$~rad at $\rho=0.06$ to $\approx0.85$~rad at $\rho=0.15$ and broadens.
% NOTE: an earlier version (RESEARCH 2026-10-01 17:20) used a window 0.5 < |Delta| < 1.8; criteria_all.py
% has no lower cut.
Even so, slow-and-tilted overcounts the bulk ($q_b\approx1.3$--$1.6$) at high $\rho$, where the wall tilt
falls to $30$--$40^\circ$.

\paragraph{Calibration against the site rule.}
\texttt{edge\_criteria.py} also determines, per run, the single threshold on $D$ (scanned over $200$ values
in $[0.5,30]$) and on $|\Delta|$ (over $[0.05,\pi]$) that maximises the agreement (fraction of $+1/2$
occurrences classified identically) with the site rule, and the agreement of the combination. In circles
with $R\le2.15\,\la$ each criterion alone agrees with the site rule on $97$--$99\%$ of occurrences
($0.89$ at $2.87\,\la$).

\paragraph{Summary of criteria per run (\texttt{criteria\_all.py}).}
For every circle run of the main sweep with $\rho=|k|/(\pi R)\le0.151$, both windings, the script computes
$q_{\rm knee}$ (knee for $k>0$, peak-plateau for $k<0$) and its position, $q_{\rm fixed}$, $q_{\rm site}$
and the number of sites, $q_{\rm dens}$ and $r_{\rm dens}$, $q_{\rm slow}$, $q_{\rm tilt}$ ($k>0$), the
$+1$ crossing $r^*/R$ and $f^*=1-(r^*/R)^2$. All inputs use $200$ shells and the 20\% transient.
The knee is the primary criterion (defined and $\approx1$ from $\sim0.6$ to $\sim2.4\,\la$ for $+k$); the
density minimum and mobility are cross-checks.

\paragraph{Layer occupancy and docking (\texttt{layer\_analysis.py}).}
Using $d_\rho$ (or $d_p$ when a passive reference exists), a docking event is a step of a same-sign track
from $s>d$ to $s\le d$, undocking the reverse; rates are events per unit time. The circulation time before
docking is the time a docking track spent in the zone $d<s\le d+2$ immediately before. In the gap test, the
angular gaps between the other layer defects at the docking frame are computed and we record whether the
docking defect enters the largest gap, compared with the random expectation $\langle 1/n_{\rm gaps}\rangle$.
The azimuthal speed $|\Delta\theta|\,\bar r/\Delta t$ of same-sign defects is binned in $2$-cell bins of $s$
(steps spanning at most 2 frames).

\subsection{Pair correlations}
\label{sec:methods-gr}

\paragraph{Lag-referenced $g_{ab}(r)$.}
In the inhomogeneous cavity the pair-distance histogram is normalised by a reference with the same one-body
density and no correlations: pairs between frame $i$ and frame $i+N_{\rm f}/2$ (cyclic), which are
decorrelated (\texttt{defect\_spatial.rdf}). For $ab\in\{++,--,+-\}$, $g_{ab}(r)$ is the ratio of the
within-frame histogram (bins of $1$ cell up to $2R$) to the cross-frame one, the latter rescaled per frame by
$\tfrac12N_a(N_a-1)/(N_aN_a')$ (same sign) or $N_aN_b/(N_aN_b')$ (opposite signs) to count pairs like the
within-frame histogram; bins with reference count $\le5$ are left blank. It is computed for all defects and
for bulk defects only ($r<R-4.2$). For static layer defects the reference contains effective self-pairs
(the passive $g$ shows the site lattice), so this estimator is meaningful only for mobile (bulk) defects.

\paragraph{Same-ring angular correlation $g_{ab}(s,\Delta\phi)$.}
The cavity is divided into rings of width $1$ in $s$; for each ring and frame, the angular separations
$\Delta\phi\in[0,\pi]$ of all pairs in the same ring are histogrammed in $90$ bins and divided by the uniform
expectation (pairs$\times\Delta\phi_{\rm bin}/\pi$); rings with fewer than $30$ pairs are blank. This needs only
within-frame pairs and no model of the radial profile; it resolves the layer lattice (peaks at
$\Delta\phi=2\pi m/(2k-2)$) but same-ring bulk pairs are too rare.

\paragraph{Ring-resolved $g_{ab}(r\,|\,s)$.}
For reference defects of type $a$ in ring $s$ (width $1$) and partners of type $b$ anywhere in the cavity,
the within-frame distance histogram (bins of $1$ cell up to $2R$, self excluded) is divided by the expected
count for an uncorrelated partner distribution with the measured radial density,
\begin{equation}
  E_{ab}(r\,|\,s) = N^{\rm ref}_a(s)\;2\pi r\,\Delta r\;\big\langle n_b\big(R-|\mathbf{r}_0+r\hat{\mathbf{e}}_\theta|\big)\big\rangle_\theta ,
\end{equation}
where $\mathbf{r}_0$ is a point at the ring's centre radius, the angular average uses $720$ directions,
$n_b(s)$ is the measured radial density ($\Delta s=0.25$) and $n_b=0$ outside the wall. Bins are shown if
$E>3$ and the ring has at least $30$ references. Seed-pooled maps (\texttt{edge\_sites.py}) sum observed and
expected counts over the seeds of each $(k,R,\zeta)$ before dividing ($E>5$); zone curves sum over rings in
three zones (edge $s<4.2$, outer bulk $4.2<s<R/2$, centre $s>R/2$) and over all $k$ at equal $R$ ($E>20$), with
a Poisson error $g/\sqrt{O}$. The thin layer peak in $n_b$ produces an artefact at the largest $r$ (partner in
the far-side layer).
% NOTE: the zone pooling uses R/2 of the first (k, R) group for all groups, i.e. it assumes equal R.

\subsection{Circulation order parameters}
\label{sec:methods-circ}

The coarse-grained velocity field (\texttt{flowfield.dat}; each frame is the average over the output
interval, every $1000$ steps $=100$ time units in the main sweep) is read per cell (cell centres at
integer$+\tfrac12$). Over the cells with $0.5<r<R-1$ (\texttt{--rcut 1}) we compute per frame
\begin{equation}
  m(t) = \big\langle \mathbf{v}\cdot\hat{\boldsymbol\phi}\big\rangle_{\rm cells},\qquad
  u(t) = \big\langle |\mathbf{v}|\big\rangle_{\rm cells},
\end{equation}
and the dimensionless order parameters
\begin{equation}
  \Phi_{\rm persist} = \frac{|\langle m\rangle_t|}{\langle u\rangle_t},\qquad
  \Phi_{\rm inst} = \frac{\langle |m|\rangle_t}{\langle u\rangle_t}.
\end{equation}
$\Phi_{\rm persist}$ measures a persistent sense of rotation; $\Phi_{\rm inst}$ is insensitive to
reversals of the spontaneously chosen sense. $\Phi_{\rm persist}\approx\Phi_{\rm inst}$ indicates a steady vortex,
$\Phi_{\rm inst}\gg\Phi_{\rm persist}$ a reversing one (e.g.\ $0.16$ vs $0.71$ in one test seed). The raw
$|\langle m\rangle_t|$ and $\langle|m|\rangle_t$ are also reported. The same script computes the bulk charge with
the fixed cut $r<R-2.5$, $P_1$, the audit fractions, and $P_1$, $\langle q_b\rangle$ restricted to audit-passing
frames (a difference would indicate that the audit filter biases the result) (\texttt{circulation.py}).

\subsection{Aggregation over seeds}
\label{sec:methods-aggregate}

\begin{itemize}
  \item \emph{Within a run}: time averages over the retained frames. Statistical errors of $\langle q_b\rangle$
        and $P_1$ are block errors with $10$ consecutive blocks (standard deviation of block means$/\sqrt{10}$;
        \texttt{layer\_analysis} truncates to equal blocks, \texttt{edge\_sites} uses near-equal splits).
  \item \emph{Parameter sets}: runs in seed directories (\texttt{.../rho*/k*/sNN}) are grouped by their parent
        directory; set values are the mean over seeds and errors the standard error,
        $\mathrm{std}_{\rm ddof=1}/\sqrt{n_{\rm seeds}}$ (\texttt{circulation.py}).
  \item \emph{Criteria maps}: per $(\mathrm{sgn}\,k,\rho,|k|)$ the bulk charge of each criterion is the mean over the
        seeds for which it is defined; a point is marked undefined if no seed has it (e.g.\ no knee) and flagged if
        it is defined in fewer than $2/3$ of the seeds (\texttt{plots\_criteria.py}, from \texttt{criteria\_all.csv}).
  \item \emph{Profiles}: seed-averaged $n_\pm(c)$ and $Q(c)$ are arithmetic means of the per-seed profiles on the
        common shell grid. Pair correlations are pooled by summing observed and expected counts before dividing
        (Sec.~\ref{sec:methods-gr}).
  \item \emph{Main-sweep table} (\texttt{main\_sweep\_analysis.py}): one row per run with $R$, $k$, $\rho=|k|/(\pi R)$,
        seed, $\zeta$, number of frames, $Q_{\rm knee}$, $c_{\rm knee}$ and contrast, $c^*/R$, $d^*=R(1-c^*/R)$,
        $f^*=1-(c^*/R)^2$, number of sites, $q_{\rm site}$, $P_1^{\rm site}$, and the wall Fourier quantities
        ($m^*$, $S(m^*)$, $f_{\rm static}$, $N_{\rm wall}$ from \texttt{wall\_fourier.py}, $+1/2$ within $4$ cells of
        the wall). The analysis is restartable (runs already in the table are skipped).
\end{itemize}
Lengths are normalised by the active length measured in periodic boxes at the same activity,
$\la(\zeta=0.07)=22.2\pm0.3$ (value $22.21$ used in the scripts), and densities by $\rho=|k|/(\pi R)$ for circles or
$2|k|/P$ for general shapes.
